\chapter[Iterations on MOC Fixes: Literature, Decisions, and Evidence]{Iterations on MOC Fixes:\\Literature, Decisions, and Evidence}
\label{ch:moc-iterations}

\lead{Development record: 19--20 November 2026.}
This chapter records the findings, implementation changes, failed attempts,
and validation results for the method-of-characteristics (MOC) and Rao
nozzle solver. It supplements the dated snapshots in earlier chapters.
The entries explain which gates failed, what caused those failures, what
changed, and how each change was checked. A passing published contour or
integral benchmark is distinct from a passing local flow-field audit, and
both are distinct from validation against measured hardware.

\textbf{Final bounded outcome.} The repaired, explicitly configured Nozzle-B
case passes all internal gates and the unchanged 32/48-to-48/64-point
common-coordinate refinement policy. Maximum field-angle difference is
0.0228912 degrees against a 0.05-degree limit, with 95.027-percent minimum
coverage. The accepted result requires exact D/E attachment, accurate seed
controls and sufficient kernel capacity; earlier coarse failures remain in
the record. This is a numerical verification result for one inviscid,
constant-$\gamma$ model. The original KL table, broad-domain accuracy and
physical validation remain unresolved. Iterations 18--19 contain the final
configuration, completed tests, independent review and reproduction command.

\section{Baseline and evidence}
The initial investigation compared the existing knowledge graph, original
source PDFs, implementation, and tests. Source equations were checked
against the original documents where available. Differences between valid
conventions and intentional model restrictions were distinguished from
implementation defects.

The starting checkout is \code{d636d854} (the complete identifier and source hashes
are in \code{analysis/baseline_manifest.json}). The working
tree already contained extensive edits, including a digitized SP-125 TOP
chart. Existing changes were preserved. Working evidence was stored below
\code{builds/moc_audit_2026-11-19/}. Durable copies of the research reports,
completed test checkpoints and failed numerical attempts are retained in
\code{docs/subsystems/S06-rao-moc-bvp-contour/audit_2026-11-19/evidence/};
short evidence paths in this chapter are relative to that directory, and
its manifest records payload hashes. The consolidated source inventory,
literature inventory, graph assessment, traceability matrix and dependency-ordered
backlog are in that same audit directory. A retained failed attempt remains
failed when later work succeeds.

The preceding completed focused run was
\code{.venv-jax/bin/python -m pytest -q tests/test_moc.py tests/test_moc_bounded_net.py tests/test_characteristic_pairing.py tests/test_nasa_kernel_march_parity.py tests/test_nasa_port.py}:
77 passed and 2 failed in 77.06 seconds. The failures were the chart-bell
coalescence-location/refinement expectation and a strict-wall-intersection fixture
whose hardcoded 30-degree arc no longer joins the current chart bell. These are
recorded separately from the nonlinear BVP and physical-topology gates.

\section{Implementation inventory and validation checks}
\begin{longtable}{@{}P{0.31\linewidth}P{0.62\linewidth}@{}}
\toprule\textbf{Module or family} & \textbf{Responsibility and audit boundary}\\\midrule
\endhead
\code{raosim/moc.py} & Local point processes and forward-analysis entry point; deprecated shrinking-front march retained for historical regression.\\
\code{raosim/moc_bounded.py} & Reflected, bounded forward net for a prescribed wall; exit sampling; containment, coalescence and mass diagnostics.\\
\code{raosim/nasa_moc.py} & Transonic starting line, throat kernel, DE integration and D placement, BDE construction, streamline wall, full control-surface mass and thrust.\\
\code{raosim/rao_variational.py} & Nonlinear control-surface residuals, construction, independent topology audit, numerical policy, reliability and convergence reporting.\\
\code{raosim/rao_residuals.py} & Characteristic compatibility, Mach-line direction and wall-tangency residuals.\\
\code{raosim/moc_topology.py} & Explicit B, D, E, BD, DE and wall curves; seam and mass-closure reports.\\
\code{raosim/jax/} & Differentiable residual assembly and throat-kernel port; parity with host implementation is a separate verification requirement.\\
\code{raosim/benchmarks.py} & Published shape/integral comparisons; a benchmark pass does not override a failed per-run physical gate.\\
\code{raosim/rao_existence_scan.py} & Fixed-context closure roots and resolution/domain evidence; evidence cannot be transferred to an arbitrary configuration.\\
\code{raosim/mdo/contour_validation.py} & Optional post-solve comparison, not a replacement MOC contour provider inside the MDO loop.\\
\bottomrule
\end{longtable}

\section{Literature inventory and graph coverage}
A literature inventory catalogued 15 relevant source payloads and recorded their SHA-256,
manifest identifier, actual title, and reviewed page windows in
\code{literature/literature_inventory.json}. Catalogue inclusion is not full review.
The original-PDF evidence below is distinguished from supporting literature which
was only inventoried. A separate knowledge-graph survey checked
the original Rao and NASA pages; its coverage and equation inventories are in
\code{knowledge_graph/coverage_inventory.json} and \code{knowledge_graph/rao_equation_inventory.json}.

\begin{longtable}{@{}P{0.28\linewidth}P{0.22\linewidth}P{0.43\linewidth}@{}}
\toprule\textbf{Source / identity} & \textbf{PDF pages checked} & \textbf{Load-bearing scope or limitation}\\\midrule
\endhead
Rao, 1958; DOI 10.2514/8.7324 & 2--6 (printed 378--382) & Full CDE fluxes, continuous D data, DE characteristic/stationarity, construction, Nozzle-B tables; equations on pages 3--4 visually inspected.\\
Rice / JHU, 2003; NTRS 20030067852 & 7--10, 16--21 & Characteristic topology, starting line, mass check, separate resolution controls, convergence warnings. Printed compatibility notation conflicts with Anderson; see below.\\
Anderson, Modern Compressible Flow, third edition; hash-named held PDF & 94--98, 419--423; especially 422 / printed 406 & Independent governing-equation verification, Eqs.~11.52--11.53; critical-Mach normalization, Eq.~3.37.\\
MIT 16.512, Martinez-Sanchez, held lecture 4--5 PDF & 1--10 & Independent characteristic derivation; family names differ, so map by geometric slope.\\
Young, Auburn thesis, 2012 & 36--40, 62--63, 94 & Geometric viability, Eq.~3.10; reacting-flow source terms are additional physics, not merely variable gamma.\\
NASA SP-8120; NTRS 19770009165 & 27, 29, 79--81 & Optimum versus approximate contours, gas-property and viscous corrections; intentional applicability limits.\\
NASA RP-1104; NTRS 19830016278 & 1--2, 5--7 & Perfect-bell family and applicability; curve ordinates were not digitized for this audit.\\
Rao, 1961; DOI 10.2514/8.5837 & 1, 3--4 & Contour approximations and applicability; not an original TOP numerical table.\\
Cuffel, Back, Massier, 1969; DOI 10.2514/3.5349 & 1--2 & Independent transonic measurements and KL comparison; no disputed KL coefficient table.\\
Frey and Hagemann, 2000; DOI 10.2514/2.5593 & 1--4, 6 & Internal shocks and separated flow; does not certify a numerical fold's physical location.\\
Ostlund, KTH thesis, 2002 & 1, 10, 122--123, 129 & Metadata and targeted coefficient search; original KL table not located.\\
Rao, Beck and Booth, 1999; AIAA 99-2584 & 3--4, targeted follow-up & Original Eqs.~(2)--(6) and the positive-radial-increment requirement checked against the existing typed validity implementation.\\
Hagemann, 1998; Rocket Nozzle Geometries; Schomberg, 2018 & Catalogue only in this independent pass & Retained as potential follow-up sources; catalogue inclusion does not certify their formulas.\\
\bottomrule
\end{longtable}

\subsection{What the knowledge graph establishes}
The synchronized corpus contained 296 on-disk PDFs and 262 canonical payloads.
The graph contained nine stored Rao-1958 equations and five claims, but no extracted
equations or claims for the NASA/JHU report, Young thesis, Ostlund thesis, or
Rao-1999 paper. The MOC model has no current independent source-review result;
all 25 inspected MOC/NASA software scientific comparisons were marked
\code{NOT_RECHECKED}. The graph's legacy ``optimizer coupled'' label is a
keyword heuristic with execution reach unknown. It is not evidence that MOC
runs inside the MDO optimization loop. Source metadata also has discrepancies:
for example, the held Anderson textbook carries a manifest identifier unrelated
to its actual title. The inventory preserves inspected and stored identities.
No graph stores were silently rewritten during this audit.

\subsection{Reconciled equations and preserved disagreements}
Anderson PDF page 422, Eqs.~11.52--11.53, supports the implemented relations
\begin{align}
C^- &: \quad d(\theta+\nu)=S\,ds, &
C^+ &: \quad d(\theta-\nu)=-S\,ds, &
S &= \frac{\sin\theta\sin\mu}{r}.
\end{align}
Here $\nu$ is the Prandtl--Meyer function and $\mu=\sin^{-1}(1/M)$.
The Rice report visibly prints $\mu$ in its compatibility expressions on
PDF pages 8--9 while also defining it as Mach angle. That discrepancy remains
documented; the independently derived Anderson equations resolve the
implementation requirement. No sign rewrite is justified.

Rao's parameter in PDF page 4, Eq.~(17), is
$M/\sqrt{2+(\gamma-1)M^2}$; Anderson's critical Mach $V/a^*$ is
$\sqrt{\gamma+1}$ times this value. At fixed $\gamma$, that constant can be
absorbed into the stationarity constant. This is a valid normalization
difference. A separate unresolved issue is the transonic kernel's use of
$\sqrt{U^2+V^2}$ as local Mach while its field names say $U=u/a^*$ and
$V=v/a^*$. The original KL normalization must be established before changing
the numerical branch.

The local KL coefficient docstring claims exact checks against Kliegel and
Levine's 1969 Table 1, but that original paper was not found in the corpus.
Targeted publisher retrieval returned inaccessible/403 responses;
NASA searches established bibliographic identity, not the coefficient table.
Power-series patterns and agreement with a historical fixture cannot verify
those exact coefficient edits. They remain an explicit research question,
not a reason to invent a replacement coefficient.

The free-exit and fixed-end variational problems also differ legitimately.
Rao's PDF page 3, Eq.~(6), retains an endpoint variation because only
length is prescribed and the exit radius is free; Eqs.~(10) and (14)
follow from that free variation. When both throat radius and expansion
ratio are fixed, $\delta r_E=0$ and that endpoint condition does not follow.
Leaving it inactive in the repository's fixed-end residual set is not
automatically a missing equation. A separate derivation showed that the
second invariant in Eq.~(13) follows from Eq.~(12), characteristic geometry
and exact compatibility on a smooth constant-$\gamma$ DE. Adding it as
another weighted residual is not justified solely because it appears
separately in the paper. Its discrete drift can still diagnose inconsistency.

Rice PDF page 19 calls its Mach-1.5 cap and starting-line shape arbitrary
numerical conventions; its 2-percent mass check and 0.25--0.5-degree nominal
turning increments are implementation policies. Neither source assigns the
repository's finite acceptance thresholds. Cuffel et al., PDF page 2, reports
a KL/ideal-one-dimensional mass coefficient of 0.982 versus 0.985 measured
for its particular throat. Thus conservation relative to the modeled start
line must remain separate from ideal-throat normalization.

\section{Traceability and prioritized backlog}
The classification describes the audited baseline, while the status records
what this cycle established or repaired. ``Missing'' can describe an intentional
model restriction or absent validation, not necessarily a code defect. Detailed
task commands, hashes and outcomes accompany the machine-readable matrix in
\code{docs/subsystems/S06-rao-moc-bvp-contour/audit_2026-11-19/traceability.json}.
Priority was correctness and refusal of invalid results first, then boundary
and conditioning repairs, then trustworthy numerical acceptance and reporting.
Unverified coefficients and additional physics were not filled in speculatively.
\begingroup\footnotesize
\begin{longtable}{@{}P{0.12\linewidth}P{0.40\linewidth}P{0.40\linewidth}@{}}
\toprule\textbf{ID / baseline class} & \textbf{Finding and implementation} & \textbf{Source, action and current status}\\\midrule
\endhead
\textbf{MOC-01}\par Correct & Characteristic-family pairing and axisymmetric source\par \code{raosim/rao_residuals.py:residual_Cplus_axisym}\par \code{raosim/moc.py:solve_interior_point} & \textbf{Evidence.} Anderson PDF422/printed406 Eqs11.52-53; Rao1958 PDF3 Eq15\par \textbf{Action.} Preserve signs; retain reference reconstruction tests\par \textbf{Status.} preserved\\ \addlinespace[5pt]
\textbf{MOC-02}\par Incorrect & Configured downstream radius ignored in preferred seed\par \code{raosim/rao_variational.py:_initial_ce_from_kernel}\par \code{raosim/nasa_moc.py:set_theta_b} & \textbf{Evidence.} Existing configuration contract; Rao1958 PDF3-4 Rd=.45Rt\par \textbf{Action.} T02b: propagate Rd; preserve default and length convention\par \textbf{Status.} implemented focused validated\\ \addlinespace[5pt]
\textbf{MOC-03}\par Incorrect & Changed fixture geometry invalidates historical regression\par \code{tests/test_moc_bounded_net.py:_bezier_wall} & \textbf{Evidence.} Existing wall seam and reference-regression contracts\par \textbf{Action.} T01: pin original angles and add current-chart fail-closed regression\par \textbf{Status.} implemented focused validated\\ \addlinespace[5pt]
\textbf{MOC-04}\par Incomplete & D boundary data treated as soft duplicated unknowns\par \code{raosim/rao_variational.py:_ce_geometry_residuals} & \textbf{Evidence.} Rao1958 PDF3-4 continuity and Eqs17-19; Rice PDF9-10\par \textbf{Action.} T03: opt-in reduced-vector exact boundary attachment\par \textbf{Status.} implemented and focused validated\\ \addlinespace[5pt]
\textbf{MOC-05}\par Incomplete & Backward BDE compatibility underresolved on coarse mesh\par \code{raosim/nasa_moc.py:_calc_bde_back_point}\par \code{raosim/rao_variational.py:_bde_measured_topology_report} & \textbf{Evidence.} Anderson PDF422 Eqs11.52-53; Rice PDF21 numerical refinement\par \textbf{Action.} Measure edge-specific residuals and vary relevant resolutions without loosening thresholds\par \textbf{Status.} local compatibility and documented refinement pair pass\\ \addlinespace[5pt]
\textbf{MOC-06}\par Incorrect & Warning conflates optimizer termination with residual exceedance\par \code{raosim/rao_variational.py:solve_rao_bvp} & \textbf{Evidence.} Existing gate bvp\_ok includes optimizer\_success plus numerical residuals\par \textbf{Action.} T04: structured gate report with actual separate failure reasons\par \textbf{Status.} implemented and focused validated\\ \addlinespace[5pt]
\textbf{MOC-07}\par Incomplete & NozzleB comparison omits differing throat and start-line assumptions\par \code{raosim/benchmark_data/cases/rao1958_nozzle_b.json}\par \code{raosim/benchmarks.py:_build_contour} & \textbf{Evidence.} Rao1958 PDF3-5; original Rd=.45Rt and Sauer start versus .382/KL model\par \textbf{Action.} T05: explicit options and applicability metadata; retain reference-profile distinction\par \textbf{Status.} implemented and focused validated\\ \addlinespace[5pt]
\textbf{MOC-08}\par Uncertain & Unverified KL polynomial changes and primary-source claims\par \code{raosim/transonic_kernel.py:_kl_axi_coeffs} & \textbf{Evidence.} Original KL1969 absent locally; publisher access inaccessible; NASA report only high-level\par \textbf{Action.} T06 correct provenance wording only; block coefficient changes pending primary table\par \textbf{Status.} provenance corrected primary math still unverified\\ \addlinespace[5pt]
\textbf{MOC-09}\par Uncertain & Transonic U/a-star labels versus returned Mach convention\par \code{raosim/transonic_kernel.py:kliegel_levine} & \textbf{Evidence.} Anderson PDF98 Eq3.37; KL original normalization not verified\par \textbf{Action.} Preserve branch behavior; require primary normalization evidence and independent benchmark\par \textbf{Status.} research blocked primary source\\ \addlinespace[5pt]
\textbf{MOC-10}\par Correct & Rao versus Anderson M-star constant scaling\par \code{raosim/gas_dynamics.py:mstar_from_M}\par \code{raosim/rao_variational.py:rao_stationarity_residual} & \textbf{Evidence.} Rao1958 PDF4 Eq17; Anderson PDF94 and98\par \textbf{Action.} Constant factor absorbable in stationarity C at fixed gamma; no rewrite\par \textbf{Status.} preserved\\ \addlinespace[5pt]
\textbf{MOC-11}\par Correct & Free-exit transversality omitted for fixed-area problem\par \code{raosim/rao_variational.py:CHARACTERISTIC_RAO_RESIDUAL_BLOCKS} & \textbf{Evidence.} Rao1958 PDF3 Eq6 endpoint variation, Eq10 and14\par \textbf{Action.} Retain fixed-end versus free-exit distinction; no unjustified residual addition\par \textbf{Status.} preserved\\ \addlinespace[5pt]
\textbf{MOC-12}\par Correct & Eq13 not a separate independent continuum residual\par \code{raosim/rao_variational.py:_rao_algebraic_stationarity_residuals} & \textbf{Evidence.} Rao1958 PDF3 Eq12,13,15; differential-identity derivation, finding TB11\par \textbf{Action.} No redundant equation; finite invariant drift can be diagnostic\par \textbf{Status.} preserved\\ \addlinespace[5pt]
\textbf{MOC-13}\par Missing & Detailed source-equation and current-review coverage in graph\par \code{knowledge/equations.json}\par \code{knowledge/models.json}\par \code{knowledge/software_links.json} & \textbf{Evidence.} knowledge-graph coverage inventory; original sources checked\par \textbf{Action.} Deliver coverage report and evidence; do not silently rewrite unrelated graph stores\par \textbf{Status.} documented gap\\ \addlinespace[5pt]
\textbf{MOC-14}\par Uncertain & Exact-axis source limit versus guards\par \code{raosim/rao_residuals.py:_source_axisym}\par \code{raosim/moc.py:solve_axis_point} & \textbf{Evidence.} Anderson PDF422 plus explicitly assumed smooth-axis symmetry expansion\par \textbf{Action.} Current midpoint callers not established defective; keep numerical evidence and research caveat\par \textbf{Status.} no confirmed active defect\\ \addlinespace[5pt]
\textbf{MOC-15}\par Missing & General accuracy envelope and independent second full-flow oracle\par \code{raosim/model_registry.py:nozzle.moc}\par \code{tests/test_moc_bounded_net.py} & \textbf{Evidence.} One NASA historical reference plus NozzleB shape/integral is limited evidence\par \textbf{Action.} Report exact tested configurations; common-coordinate refinement and new oracle remain separate work\par \textbf{Status.} validation gap\\ \addlinespace[5pt]
\textbf{MOC-16}\par Missing & Reacting/variable-gas, shock-fitting, viscous and full JAX forward-net extensions\par \code{raosim/moc_bounded.py}\par \code{raosim/jax/moc_kernel.py} & \textbf{Evidence.} Young PDF94 additional reactive source terms; SP8120 PDF79-81; Frey2000\par \textbf{Action.} Intentional restricted model; outside immediate confirmed-gate repair, no speculative physics\par \textbf{Status.} outside current repair scope\\ \addlinespace[5pt]
\textbf{MOC-17}\par Incorrect & Unit-process iteration exhaustion accepted as a valid point\par \code{raosim/nasa_moc.py:_calc_arc_wall_point_raw}\par \code{_calc_bde_back_point} & \textbf{Evidence.} NASA C++ CalcArcWallPoint lines878-925 and CalcBDERegion3331-3402; conCrit line62; source hash in the unit-process convergence contract\par \textbf{Action.} T07 restore failure-on-exhaustion and exact-zero-safe local iteration checks without changing tolerance\par \textbf{Status.} implemented and focused validated\\ \addlinespace[5pt]
\textbf{MOC-18}\par Incorrect & Refinement identity omitted boundary/model differences\par \code{raosim/rao_variational.py:compare_rao_solutions_common_coordinates} & \textbf{Evidence.} Existing comparison same-physical-model contract; Rao shared-boundary construction; Rice starting-line/curvature conventions\par \textbf{Action.} T08 require matching curvature, starting-line, attachment and residual recipe provenance; missing identity fails closed\par \textbf{Status.} implemented and focused validated\\ \addlinespace[5pt]
\textbf{MOC-19}\par Incorrect & Inactive multiplier unknowns create exactly zero Jacobian columns\par \code{raosim/rao_variational.py:_exact_d_free_indices} & \textbf{Evidence.} Configured residual dependency contract; AD rank29of31 at real NozzleB seed\par \textbf{Action.} T09 remove only provably inactive lambda2/lambda3 in opt-in mode; after reduction rank29of29\par \textbf{Status.} implemented and focused validated\\ \addlinespace[5pt]
\textbf{MOC-20}\par Missing & Reduced-coordinate derivative chain absent from sensitivity API\par \code{raosim/jax/sensitivities.py:rao_sensitivities} & \textbf{Evidence.} Mathematical chain rule for newly constrained D boundary; independent code review\par \textbf{Action.} T10 explicitly refuse unsupported exact-D sensitivities before solve\par \textbf{Status.} implemented explicit refusal\\ \addlinespace[5pt]
\textbf{MOC-21}\par Incorrect & CLI validity gate used legacy Eq6 proxy rather than typed Eq5 result\par \code{raosim/run_nozzle.py:_rao_radial_validity_gate} & \textbf{Evidence.} Rao1999 PDF3 Eqs2-3 and PDF4 Eqs5-6; existing typed validity contract independently source-checked\par \textbf{Action.} T11 use typed available+valid; missing refuses; expose actual solver gate report\par \textbf{Status.} implemented and focused validated\\ \addlinespace[5pt]
\textbf{MOC-22}\par Incomplete & Single internal-gate pass does not establish common-coordinate resolution acceptance\par \code{scripts/audit_moc_reference_profile.py}\par \code{raosim/rao_variational.py:compare_rao_solutions_common_coordinates} & \textbf{Evidence.} Rice PDF21 convergence guidance; existing common-coordinate policy\par \textbf{Action.} T12/T13/T18/T19 retain failed coarser pairs and check one further pair under unchanged field policy\par \textbf{Status.} documented final pair passes coarser failures retained\\ \addlinespace[5pt]
\textbf{MOC-23}\par Incorrect & Truncated throat kernel accepted as a converged-angle seed after row cap\par \code{raosim/nasa_moc.py:set_theta_b}\par \code{raosim/rao_variational.py:_initial_ce_from_kernel}\par \code{raosim/jax/theta_b_solve.py:make_residual_theta_b} & \textbf{Evidence.} Rice PDF7-10 full kernel/known BD construction; existing reached\_wall contract; direct fixed-angle500vs1000 rows test\par \textbf{Action.} T15 expose capacity, reject incomplete kernels before DE or BVP use, preserve explicit computational versus angle failures, repeat finer profile\par \textbf{Status.} implemented and focused validated\\ \addlinespace[5pt]
\textbf{MOC-24}\par Incorrect & Actual final live JAX kernel completeness flag ignored\par \code{raosim/jax/theta_b_solve.py:_require_final_live_kernel}\par \code{make_residual_theta_b} & \textbf{Evidence.} Existing reached\_wall contract; Rice PDF7-10 complete throat-kernel B-D-E construction; independent review of static allocation estimate versus actual raw steps\par \textbf{Action.} T17 recompute final live march and refuse incomplete result; no claim that trial iterates are qualified\par \textbf{Status.} implemented and focused validated\\ \addlinespace[5pt]
\textbf{MOC-25}\par Incomplete & Fixed eight-iteration frozen-angle seed insufficient for tighter refinement checks; termination accuracy not exposed\par \code{raosim/nasa_moc.py:set_theta_b}\par \code{raosim/rao_variational.py:_initial_ce_from_kernel}\par \code{RaoSolverConfig} & \textbf{Evidence.} Existing fixed-end seed residual contract; controlled 8-versus25-versus40 iteration experiments; Rice PDF9-10 and21 numerical convergence conventions\par \textbf{Action.} T16 expose budget/tolerance without changing defaults, report actual residual and termination reason, then rerun comparison\par \textbf{Status.} implemented focused validated honest seed nonconvergence report retained\\ \addlinespace[5pt]
\textbf{MOC-26}\par Incomplete & Soft fixed-exit endpoint can stop short of prescribed L despite internal gate pass\par \code{raosim/rao_variational.py:_exact_exit_boundary_data}\par \code{_expand_exact_d_unknowns}\par \code{raosim/jax/assembly.py:expand_exact_d_unknowns} & \textbf{Evidence.} Rice PDF9-10 prescribed-endpoint construction; project derivation of existing midpoint C+ geometry; independent positive-branch and chain-rule review\par \textbf{Action.} T18 opt-in exact endpoint elimination with domain penalties/final refusal, unchanged export and physical criteria\par \textbf{Status.} implemented and validated static profile and refinement pair\\ \addlinespace[5pt]
\bottomrule
\end{longtable}
\endgroup


\section{Iteration 01: Repair the identity of the test geometry}
\subsection{Issue and intended change}
The TOP angle-table update changed the default bell at expansion ratio 10 and
80-percent length from 30/15.5 degrees to 25.74/10.49 degrees. One test still
constructed a 30-degree throat arc, and another asserted the older wall's fold
location for the changed wall. The fixture repair separated the historical
explicitly specified wall regression from current-chart behavior.
The numerical solver and original historical location/refinement tolerances
were retained. Current-chart failure detection was recorded separately from
resolution convergence.

\subsection{Change and validation}
Only \code{tests/test_moc_bounded_net.py} changed: the wall helper accepts
explicit angles; the historical fold-location regression now specifies
30/15.5 degrees; a separate current-chart test requires honest incomplete status,
a fold inside the nozzle, and zero containment violations; and the strict-wall
test builds its arc from the actual contour. No production algorithm or
numerical threshold changed.

The four-grid diagnostic at $(n,\Delta\theta)=(24,0.5^\circ)$,
$(48,0.25^\circ)$, $(96,0.125^\circ)$, and $(144,0.0625^\circ)$ placed the
historical-wall fold at 22.629, 21.832, 20.922, and 20.973 mm. The current-chart
wall gave 55.269, 51.569, 51.788, and 50.727 mm. These are different nozzle
geometries. The second sequence is not accepted here as a resolution-convergence
demonstration. Full inputs and hashes are in
\code{validation/t01_coalescence_probe.json}.

\textbf{Completed validation:}
\code{.venv-jax/bin/python -m pytest -q tests/test_moc_bounded_net.py tests/test_nasa_kernel_march_parity.py tests/test_nasa_port.py}
returned exit 0: 31 passed in 22.58 seconds. The command, final output, and
validation record are in \code{validation/t01_result.json} and
\code{validation/t01_pytest.log} below the audit evidence directory. These results
validate the fixture repair and reference regressions, not the full BVP.

\section{Iteration 02: Explain the nonlinear and topology failures separately}
\subsection{Measured baseline}
The exact Nozzle-B benchmark configuration was rerun with the NumPy backend,
10 control points, 10 starting points, 100 function evaluations, and a
\code{residual_tol} of 0.005. The largest weighted residual was 0.003979,
dominated by regularization; mass and length residuals were $-0.001478$ and
0.000508. Nevertheless the optimizer did not declare successful termination.
This is distinct from the physical-mesh failures:
\begin{itemize}
\item The DE starting radius misses kernel point D by $0.00121693R_t$,
while the seam policy permits $0.00001R_t$. The current least-squares
formulation treats the boundary values as soft residuals.
\item The largest independently measured C+ compatibility error is
0.012274 degrees against a 0.005-degree policy. C- reaches 0.001985 degrees.
The optimizer's C+ residual is weighted by 0.05, so its reported scaled norm
does not directly certify the unweighted flow audit.
\item Crossing, wall-tangency, Mach-line direction, cell orientation, axis,
axial mass, wall geometry, and mass-closure checks pass in this configuration.
The failed seam also causes physical-mesh completeness to fail. The auxiliary
caustic is downstream of the exit; that fact alone is not the failed physical gate.
\end{itemize}
The machine record is \code{analysis/nozzle_b_baseline.json}. A 24-point,
unsolved reference seed at expansion ratio 10 has exact D attachment and a
complete physical strip but still fails measured C+ compatibility (maximum
0.022761 degrees). That seed is not a converged solution and is not a
counterexample to the independent gate. It is recorded in
\code{analysis/reference_seed_baseline.json}.

\subsection{Literature requirement and proposed investigation}
Rao (1958), PDF page 3 / printed page 379, states continuity of interior Mach
number and flow angle and continuity of the control-surface inclination; page 4,
Eqs.~(17)--(19), locates D on the kernel characteristic and constructs DE from
that state. Rice/JHU, PDF pages 9--10, describes the same boundary construction.
The shared D state is a physical boundary requirement. Enforcing it by variable
elimination rather than penalties would be an implementation choice, and must
preserve equivalent host and JAX formulations.

The diagnostic probes varied solver budget, residual weight, control resolution and
diagnostic regularization to locate the dominant physical edges and distinguish
budget exhaustion, boundary tradeoffs, and discretization. Acceptance tolerances
were held fixed. The completed probes show that increasing the 0.05-weight
evaluation budget from 100 to 1,000 reduces the D gap to $0.000341R_t$ and
C+ maximum to 0.00812 degrees; both still fail. Raising the physics weight to
1 does not solve the interior error: its worst edge is a backward-marched
near-wall cell, not the prescribed DE boundary. Increasing control resolution
from 10 to 16 points in a diagnostic no-regularization run reduces that
compatibility maximum from 0.01534 to 0.003729 degrees, while the D seam
still fails. Removing regularization was a diagnostic ablation only and is
not adopted as a default change. These results justify repairing the boundary
parameterization and evaluating physical-mesh resolution separately, rather
than merely increasing iteration counts or penalty weights.

\section{Iteration 03: Restore downstream throat-radius propagation}
\subsection{Confirmed defect and specification}
The worked Rao-1958 construction uses $R_u=1.5R_t$ and
$R_d=0.45R_t$ (PDF pages 3--4). The current Nozzle-B benchmark instead
inherits the repository's $0.382R_t$ downstream default and a different
starting-line model. A probe requesting $0.45R_t$ unexpectedly produced
bit-identical results to $0.382R_t$. This was traced to the preferred
\code{_initial_ce_from_kernel -> set_theta_b} path: the first function
computed the requested radius but the second hardcoded $0.382R_t$.
Fallback kernel construction already honored the requested value, so the
same configuration could describe different geometry depending on seed success.

An optional keyword \code{Rd} was added to \code{set_theta_b}, with
the historical default preserved and non-finite/non-positive/bool overrides
refused. The preferred seed now forwards the requested absolute radius.
The existing sharp-corner length-percentage convention is unchanged. This
repair makes the input effective; it does not claim exact Rao-source
reproduction, because the transonic construction still differs.

\subsection{Change and validation}
The change updated \code{raosim/nasa_moc.py} and
\code{raosim/rao_variational.py}, and added
\code{tests/test_moc_boundary_contracts.py}. Real preferred-seed tests
inspect the actual kernel radius and B coordinates for 0.382 and 0.45,
rather than only checking that a keyword was passed.

\textbf{Completed validation:}
\code{.venv-jax/bin/python -m pytest -q tests/test_moc_boundary_contracts.py tests/test_moc_topology.py tests/test_nasa_kernel_march_parity.py tests/test_nasa_port.py}
returned exit 0: 29 passed in 32.66 seconds, recorded in
\code{validation/t02b_pytest.log}. Default source-port controls continue to pass.

\section{Iteration 04: Impose shared D boundary data exactly}
\subsection{Boundary parameterization}
The opt-in \code{exact_d_attachment} mode removes the duplicated
starting-radius unknown and, when their existing pin flags are enabled,
the starting Mach and angle unknowns from the optimization vector. At
every residual evaluation the reduced vector is expanded into the existing
packed layout using the current arc-fraction-selected kernel point D.
Both position and selected flow-state components then represent the same
point to floating precision; the public packed representation and residual
ordering remain compatible. The JAX expansion must differentiate through
D's current arc fraction, rather than treating those data as constants.

This is an implementation choice enforcing Rao's shared boundary condition,
not a new source equation. It must preserve diagnostic pin flags and
explicitly refuse the unimplemented live-kernel-angle combination.
Acceptance requires host/JAX value and derivative checks, actual D-state
identity, preserved default behavior, and independent physical-gate results.
It does not guarantee the remaining backward mesh is resolved.

\subsection{Status}
The implementation added a host reduced-vector wrapper and matching JAX expansion;
the public packed residual layout is unchanged. Unsupported combinations
with a live solved kernel angle are explicitly rejected. Review checked the
reduced/full mappings, derivative dependence on the D arc fraction, and
diagnostic pin-flag behavior.

The focused boundary plus existing assembly-parity suite completed with
28 passed in 25.40 seconds (exit 0), including finite-difference versus
automatic-differentiation checks. The real NumPy Nozzle-B regression has
identical D radius, Mach and angle and zero geometric seam error, without
changing the $10^{-5}R_t$ seam policy. Genuine JAX optimizer execution is
recorded separately in \code{validation/t03_jax_runtime.json}; mocked wiring
tests are not substituted for that integration run.

Numerical probes with exact attachment at $(n_c,n_k,\Delta\theta)=$
$(10,10,0.5^\circ)$, $(20,24,0.25^\circ)$ and
$(40,24,0.125^\circ)$ all preserve the exact seam. Their C+ maxima are
approximately 0.01046, 0.00928 and 0.01189 degrees. Each used only
100 evaluations and unfinished optimization; this is not a clean
refinement-convergence study. The repair removes a boundary inconsistency;
it does not establish resolution of the backward characteristic strip.

\section{Iteration 05: Report the actual failing gates}
The existing BVP verdict includes successful optimizer termination as well
as assembled-residual, mass and length criteria. Previously, a generic
warning described every failure as residual nonconvergence, including the
Nozzle-B case whose residual was below tolerance but whose optimization
budget had ended. The change added \code{construction_diagnostics['gate_report']}
with separate optimizer, weighted-residual, mass, length, topology, validity,
and promotion entries. It records compatibility in degrees, seam distances
in meters, limits, failed required gates, and missing evidence. Report-only
neighbor increments are excluded from the required-gate list.

Optimizer counters distinguish NumPy function evaluations from JAX LM steps.
The requested \code{max_nfev} is not falsely described as an identical budget
across backends: JAX has a minimum per-stage step budget and continuation
stages. Warning messages now distinguish unsuccessful termination from
actual residual exceedance. Acceptance conditions and thresholds are unchanged.
The focused boundary/reporting suite completed with 21 passed in 13.04
seconds (exit 0), recorded in \code{validation/t04_pytest.log}.

\section{Iteration 06: Restore unit-process convergence failure handling}
Comparison of the port with the original, hash-identified
\code{MOC_GridCalc_BDE.cpp} exposed missing failure branches. Its arc-wall loop has a 50-iteration cap
(lines 878--879) and explicitly reports nonconvergence and exits at
lines 922--925. The backward BDE point does the same at lines 3331--3335
and 3391--3394 before accepting the state at lines 3397--3402.
Both use the $10^{-10}$ source convergence criterion (line 62).
Python's unconditional last-iterate returns omit these failure branches.

The repair rejects exhausted/non-finite points through the existing
optional-result contracts. An exact finite zero-to-zero increment must count
as zero error; the current sentinel 9.9 cannot establish that. The repair
must avoid reproducing C++ zero-over-zero NaNs and must not alter coefficients
or acceptance tolerances. Deliberately insufficient iteration budgets,
exact-zero states, source regressions and affected physical gates are the
acceptance tests. A local finite iteration-change helper was added, and
the correctors return \code{None} when the arc-wall or backward-point corrector exhausts
its budget without convergence. The broader legacy error helper and
unrelated JAX unit processes were left unchanged. Finite invalid input is
not substituted for a valid state; exact unchanged zero values are handled
without NaN arithmetic. Existing callers retain their incomplete-result
handling.

The focused iteration-failure, bounded-net, source-reference and boundary
contract suite completed with 62 passed in 37.26 seconds (exit 0), recorded
in \code{validation/t07_pytest.log}. Deliberate zero/one-iteration exhaustion
does not return the former initial or last guess as a successful point.
Convergence of an individual point remains distinct from passing the
independent characteristic equations on the mesh.

\section{Iteration 07: Make benchmark assumptions explicit}
The benchmark now passes case-specific upstream and downstream throat
ratios and the exact-attachment option to the existing solver configuration.
The original Nozzle-B profile explicitly retains $R_u/R_t=1.5$ and
$R_d/R_t=0.382$ for backwards compatibility. Its JSON and Markdown reports
now carry a model-assumptions block identifying Rao's $0.45R_t$/Sauer source
construction, the modeled $0.382R_t$/project-modified-KL construction, and
\code{exact_boundary_reproduction=false}. Published shape and integral
tolerances remain unchanged. An overall benchmark pass cannot silently
become a claim of identical boundary data or a passing local flow field.
The focused benchmark/starting-line suite completed with 16 passed in
13.91 seconds (exit 0). The full
benchmark/starting-line suite reported 20 passed and two failures in
24.21 seconds. Both failures concern the separate LEA/TOP B\'ezier length
comparison: the working geometry gives $L/R_t=15.044818739$ against
$15\pm0.02$. The HEAD and modified benchmark implementations were executed
against the same current geometry; their metrics were identical. Thus the
failures predate this reporting repair. Their evidence is preserved in
\code{validation/t05_lea_baseline_comparison.json}; no benchmark tolerance or
expected status was changed to conceal them.

\section{Iteration 08: Correct unsupported coefficient provenance claims}
The transonic kernel's long header assertions of verified
KL-1969 errata were replaced with a concise, qualified description and a
pointer to this chapter. This audit does not establish that the project modifications are
wrong; it establishes that the claimed primary-source coefficient check
cannot currently be reproduced. The misleading attribution of the
repository's $0.382R_t$ default to the original Rao-1958 example was removed.

For future research, the historical substitutions being questioned are:
the AXI third-order radial-velocity mixed bracket's multiplication versus
addition, $1181$ versus $1881$, its $y^2$ versus $y^3$ power, the AXI
third-order axial-velocity divisor 34 versus 384, and the planar mixed
bracket's multiplication versus addition. These describe existing project
edits, not new substitutions made in this audit. The source-visible C++
compatibility branch and the project-modified branch remain distinct.

Executable AST comparison with docstrings stripped confirmed that this
documentation task changed no operations or numerical coefficients.
The exact hash is recorded in \code{validation/t06_result.json}. Numerical
regression testing is recorded separately from this provenance correction.

\section{Iteration 09: Require the same model in refinement comparisons}
The common-coordinate comparison checked throat radius,
expansion ratio, gamma, ambient pressure and length, but not both throat
curvatures or the starting-line model. The new exact-attachment mode was
also absent from the comparison identity. Two different boundary models
could consequently be described as coarse and fine versions of one model.

The repair records the throat ratios, starting-line method,
attachment and pin settings, formulation, and effective residual recipe
in the solution provenance and rejects mismatched or missing required
identity before evaluating refinement. Missing provenance must not be
silently filled from current defaults. Node counts and turning limits
may vary as declared refinement controls; the underlying boundary model
must remain the same. Field-error tolerances were unchanged.
Focused validation results are recorded below.

\section{Iteration 10: Remove provably inactive optimizer columns}
The automatic-differentiation Jacobian was evaluated at the real
10-point Nozzle-B seed in exact-attachment mode. It had 69 residuals and
31 free columns, but columns 30 and 31 (the full-layout locations of
$\lambda_2$ and $\lambda_3$) were exactly zero. The local rank was 29.
This follows from the configured residual recipe: stationarity and
transversality blocks which use those multipliers are inactive; the
algebraic characteristic formulation uses its own stationarity constant.
The two zero singular values therefore do not establish a physical
non-existence result or an inherent singularity of Rao's equations.

The repair removes those two inactive optimizer unknowns only
inside the opt-in reduced mode, retaining their template values in the
public full vector. If a block which uses them is active, they remain
unknowns. The default legacy solve layout stays unchanged. Acceptance
requires backend parity, correct active-block handling, and a conditioning
check on the actual reduced system. The measured seed matrix is in
\code{analysis/conditioning_probe.json}; the acceptance evidence follows below.

\section{Iteration 11: Refuse unsupported reduced-coordinate sensitivities}
Independent review found the host and JAX boundary expansions consistent
with the source requirement. It also identified an
adjacent API issue: \code{rao_sensitivities} still differentiates the
full independent-state vector. Those are full-state partial derivatives,
not the reduced optimizer's derivatives with the eliminated D-state chain.
Its reported Jacobian conditioning would likewise describe a different
system. An explicit early refusal was added for exact-attachment
mode until the sensitivity API supports the reduced coordinates. This
preserves the meaning of the existing default API without inventing a
new total-derivative capability. Validation is recorded below.

\subsection{Validation of Iterations 09--11}
The common-coordinate identity repair passed 22 focused tests in 0.42
seconds. The reduced-column and sensitivity-refusal changes were included
in the completed 67-pass combined suite recorded in Iteration 12. An
independent repeated AD probe now has a $69\times29$ Jacobian, rank 29,
and no exactly zero columns; the smallest singular value is approximately
0.01569 at that seed. Those observations are local conditioning evidence,
not a global convergence proof.

\section{Iteration 12: Align the CLI and reports with the authoritative gates}
Integration review found that the CLI still decided its validity
gate from the legacy Eq.~(6) diagnostic \code{boundary_min}, although the
solver's authoritative criterion is the typed Rao-1999 Eq.~(5) radial
increment result. The new small helper requires that typed evidence to be
available and valid; missing evidence refuses BVP export, while the TOP
chart path remains outside this MOC validity check. Tests deliberately
make Eq.~(5) and the legacy diagnostic disagree. This aligns the CLI with
the existing source-reviewed contract; it does not introduce a new threshold.
The original Rao-1999 PDF pages 3--4 were subsequently checked:
Eqs.~(2)--(3) are on page 3 and Eqs.~(5)--(6) on page 4. The implemented
signs and factors agree. Midpoint evaluation and the $10^{-9}$ numerical
boundary band remain project policies. This is a necessary local
advancing-DE condition, not a sufficient theorem of global convergence,
arbitrary fixed-exit optimality, or absence of shocks in every flow.

The structured solver-gate report is now surfaced in CLI summaries and
benchmark JSON/Markdown. The benchmark export also forwards the same
upstream and downstream radii used by the solver, rather than silently
reverting to contour-conversion defaults. The combined boundary,
derivative-parity, sensitivity, CLI-validity, benchmark and provenance
suite completed with 67 passed and one slow test deselected in 54.50
seconds (exit 0).

\section{Iteration 13: A passing Nozzle-B configuration}
\subsection{Reproducible numerical specification}
The first resolved-profile probe uses the following existing controls,
with only the new exact-attachment option added:
\begin{lstlisting}
RaoSolverConfig(
    Rt=1.0, epsilon=4.973, gamma=1.23, pa_over_p0=0.0,
    length_pct=64.0452248702277,
    throat_upstream_radius_factor=1.5,
    throat_downstream_radius_factor=0.382,
    starting_line_method="kliegel_levine",
    n_control=16, n_kernel=24, kernel_dtheta_limit_deg=0.25,
    max_nfev=1000, residual_tol=0.005,
    evaluate_moc=True, wall_method="bde", solver_backend="numpy",
    exact_d_attachment=True, physics_weight=1.0,
)
\end{lstlisting}
The existing regularization is retained. No physical gate, tolerance,
endpoint-forcing policy, or monotonic-cleanup policy was relaxed.

\subsection{Measured outcome and its limits}
The NumPy optimizer terminated successfully by its gradient criterion in
25 function evaluations. The largest assembled residual is
$4.8106\times10^{-4}$, the D seam is exactly zero, C+ compatibility
maximum is 0.0011553 degrees and C- maximum is 0.00012017 degrees,
both below the 0.005-degree policy. All eleven physical/topology gates
pass, as do the independent radial-validity and full-CDE thrust gates.
Reliability is \code{rao_variational_residual_solved};
\code{converged=True}. The measured $C_F=1.5794212$ is approximately
0.22 percent below the published 1.5829 for the source example, while the
explicit throat/start-line model differences from Iteration 07 remain.

This demonstrates that the combined boundary, conditioning and numerical
resolution repairs can satisfy the existing gates in one documented
configuration. It does not isolate one change as the sole cause, prove
global optimality, establish a general parameter domain, verify the
unavailable original KL coefficient table, or validate hardware. The later
independent solves and common-coordinate field comparisons below test the
limits of this result. A regression retains the actual all-gate pass rather
than merely checking that a finite contour was returned.

\subsection{First refinement pair: retain the failed result}
The independent finer configuration $(n_c,n_k,\Delta\theta)=$
$(24,32,0.125^\circ)$ also converges and passes all internal gates,
with $C_F=1.57956997$ and C+ maximum 0.000323787 degrees. Nevertheless
the common-coordinate comparison of the 16/24 and 24/32 solutions
fails the unchanged field policy: maximum angle difference is
0.09237 degrees (limit 0.05 degrees), and the minimum bidirectional
field coverage is 0.88298 (minimum 0.90). Mach maximum difference
0.002255, relative thrust difference $9.417\times10^{-5}$ and maximum
wall-radius difference $0.001012R_t$ pass their respective limits.

Thus both local solves passing is insufficient to call the coarser
profile resolution-accepted. The script intentionally returns exit 2
for this two-level check. The next finer solve and comparison preserve
this failed pair rather than deleting it or changing its tolerances.

The report-only comparison with all seven published Table-2 wall states
also gives useful limits: finer-grid Mach differences are approximately
$-0.029$ to $-0.004$, and flow-angle differences approximately
$-0.953$ to $+0.335$ degrees. No new acceptance threshold was invented
for these rounded, different-boundary-model data. Passing internal
compatibility does not mean every published local state has been reproduced.

\section{Iteration 14: Diagnose the finer-grid branch change}
The third numerical level, $(32,48,0.0625^\circ)$, did not continue the
same successful branch. It returned $C_F=1.59831$ and failed BVP, MOC,
and radial-validity gates; its exit-radius mismatch was approximately
$0.2421R_t$. The kernel angle changed by 2.34375 degrees. The subsequent
field comparison failed dramatically, with Mach maximum difference
0.2347, angle difference 6.801 degrees and minimum coverage 0.1121.
These data were retained as a failed attempt; no fourth-grid extrapolation
or relaxed tolerance was justified from this trend.

Seed inspection identified a concrete computational cause:
\code{build_kernel} reached its default 500-row limit and returned
\code{reached_wall=False}, but \code{set_theta_b} still constructed a
DE topology and labeled it \code{fixed_end_secant}. The supposed finer
solution therefore started from a truncated kernel, at a different angle,
with a seed exit station 6.42449 instead of the requested 2.94.
This is not evidence that correctly resolved MOC theory diverges.

\buildfig{figures/moc_audit/kernel_capacity_and_compatibility.pdf}
{Two separate measured experiments. Left: the same 48-point starting line,
0.0625-degree turning limit and prescribed angle, with only the row capacity
changed; the initial data line adds one stored row. Right: the original and
repaired Nozzle-B numerical profiles, with their explicitly different controls
and residual weights, against the unchanged local C+ threshold. The right panel
is not an isolated mesh-order experiment or a proof of field convergence.}{0.98}

The resulting change added an explicit row-capacity control, preserving
the existing default, together with refusal of incomplete kernels before
both preferred and fallback BVP seeding. Capacity exhaustion must not be misclassified
as an insufficient-expansion angle and used to move the angle bracket in
the wrong direction. A larger explicit capacity may be used on finer
meshes; accuracy tolerances remain unchanged. The direct fixed-angle
500-versus-1,000-row experiment, source changes and repeat refinement
are recorded in the next iteration.

\subsection{Broad regression checkpoint before the row-budget repair}
The 24-file MOC, reference, residual, topology, JAX, sensitivity, CLI,
transonic and benchmark run completed with 296 passed, two failed,
20 slow tests deselected, and two expected deprecated-march warnings
in 820.05 seconds. Exit status was 1. Both failures were the independently
proven preexisting LEA/TOP length/status expectations described in
Iteration 07. The complete output is
\code{analysis/pre_t15_combined_pytest.log}, explicitly the checkpoint
\emph{before} the subsequent row-capacity repair.
Source hashes for that checkpoint are retained separately from the final
source inventory. Later changes require their own affected-path tests.

\section{Iteration 15: Make kernel capacity explicit and reject partial seeds}
The fixed-angle experiment confirmed that a 500-step budget stores 501 rows
but only reaches 15.78558 degrees of the requested 27.65625-degree arc.
With a 1,000-step budget the same calculation reaches the arc end in 880
stored rows. The initial data line explains the extra stored row.

The new \code{kernel_max_rrcs} control preserves the default 500 and
is threaded through the host kernel builders, preferred/fallback BVP seeds,
topology convenience builder and live-angle host rebuild. Invalid budgets
are rejected. The low-level builder may still return an explicitly partial
diagnostic object; high-level design consumers now reject it before DE
construction or successful seed provenance. \code{IncompleteKernelError}
escapes the angle-bracketing catches, so resource exhaustion cannot be
mistaken for an insufficient expansion angle. Completion, row count and
budget are recorded in the numerical context. Increasing capacity during
refinement is permitted; treating an incomplete kernel as a refined field is not.

The JAX live-angle setup also checks its static allocation against the
declared budget. Independent review then identified that this allocation
estimate alone is not proof of actual live-march completeness; that issue
is handled separately in Iteration 17. The host-kernel repair passed
99 focused tests in 54.30 seconds (exit 0; four slow
tests deselected). These tests include NASA controls, partial-kernel
refusals, propagation, context, AD parity, and the real passing profile.

The adequate-budget finest solve now has a complete 887-row kernel,
$\theta_B=27.8515625$ degrees, $C_F=1.57931261$, and C+ maximum
0.00015848 degrees. It passes all eleven raw physical/topology gates,
BVP, Eq.~(5) and thrust checks. Promotion nevertheless refuses because
the raw wall ends near 2.93975 instead of 2.94. The field comparison
also still fails: coverage 0.90863 now passes, but Mach maximum
0.00678, angle maximum 0.23503 degrees, angle-$B$ difference
0.1953125 degrees and D-fraction difference 0.002652 exceed their
unchanged limits. The resource defect is repaired; those remaining
accuracy and endpoint issues are not concealed.

\section{Iteration 16: Separate seed-search accuracy from final closure}
The preferred seed hardcoded eight angle-search iterations, intentionally
trading seed quality for speed while the kernel angle remained frozen in
the subsequent solve. That numerical choice was tested against
the tighter field-comparison requirements. At the 24/32 grid, increasing
the existing search budget from 8 to 25 changes $\theta_B$ from
27.65625 to 27.752341628 degrees and reduces the seed's relative exit
station error from 0.003827 to $2.043\times10^{-8}$. Forty iterations
give the same selected angle because the angular bracket-width floor
has already been reached. The requested $10^{-8}$ residual criterion
is therefore not falsely labeled satisfied.

A BVP run with that independently computed angle still undershoots the
exit, near 2.93994. Thus seed accuracy and soft endpoint constraints are
separate effects. The change exposed the existing search budget and
criterion, retained defaults 8 and $10^{-6}$, and recorded actual
iterations, best residual, criterion-met flag and termination reason.
The controls and truthful termination report passed the later 114-test
focused suite, including parameter forwarding and unchanged defaults.
These checks validate the controls and report independently of whether a
finer field comparison passes.

The completed diagnostic with independently computed accurate angles on the
24/32 and 32/48 grids confirms the separation. The angle-$B$ drift shrinks to
0.006482 degrees, field Mach maximum difference is 0.0016753, and minimum
coverage is 0.922699; those limits pass. Field angle maximum difference is
still 0.0649365 degrees against 0.05, and both walls remain short of the
export interval. Relative thrust difference is approximately
$7.267\times10^{-5}$, wall difference $0.0004589R_t$, and D-fraction
difference $2.409\times10^{-5}$. The diagnostic process completed with
exit 0 because it successfully measured the trials; its retained numerical
acceptance result is \emph{FAIL}. The distinction is explicit in
\code{validation/t16_seed_accuracy/report.json}.

\section{Iteration 17: Check the actual final live JAX kernel}
The live residual consumed \code{m.bd} without checking
\code{m.reached_wall}. The static allocation estimate counts half-angle
special steps, but accepted raw steps may be smaller; that estimate does
not prove every march reaches B. A complete host seed and later complete
host rebuild also cannot prove the optimized live field was complete.
The repair recomputes the actual final live march and refuses its
result if incomplete. This does not claim all trial iterates are qualified,
and it does not add support for combining exact attachment with live-angle
optimization. Eleven focused tests passed, two slow tests were deselected,
and the run completed in 11.98 seconds (exit 0;
\code{validation/t17_pytest.log}). Independent review also checked the final
guard's placement before residual/report/return. Allocation documentation
now states the estimate's limits. These results validate the refusal
behavior, not a full-domain live-angle solve.

\section{Iteration 18: Exact fixed-exit boundary parameterization}
The fixed-end model prescribes both $(L,R_E)$, yet finite least squares
can trade a small endpoint miss for other residuals. The exporter correctly
refuses to extend a wall which stops short. Independent checks of the
following derivation confirmed consistency with the existing midpoint geometry:
\begin{align}
\Delta x &= L-x_{n-2}, & \Delta r &= R_E-r_{n-2},\\
\beta &= \operatorname{atan2}(\Delta r,\Delta x), &
\mu_i &= \sin^{-1}(1/M_i),\\
r_{n-1} &= R_E, &
\theta_{n-1} &= 2\beta-\theta_{n-2}-\mu_{n-2}-\mu_{n-1}.
\end{align}
This solves the final discrete C+ relation
$\Delta r=\Delta x\tan[(\theta_{n-2}+\theta_{n-1}+
\mu_{n-2}+\mu_{n-1})/2]$ on its downstream, outward branch.
Both endpoint values can then be eliminated from the optimizer vector,
recomputed through the entire prefix at each evaluation, and reconstructed
to arithmetic-roundoff accuracy. There is no post-solve endpoint forcing
and no additional free-exit transversality equation.
The numerical implementation uses
$\mu_i=\sin^{-1}[1/\max(M_i,1.000001)]$, identical to both existing
geometry integrators; an independent domain check still requires $M_i>1$.

This is equivalent to the intended exact constraints, \emph{not} to the old
finite soft-penalty objective: restricting those tradeoffs is deliberate.
Rice PDF pages 9--10 supplies the prescribed-endpoint construction; the
algebraic elimination is an implementation derivation. The new mode is
opt-in and requires the static exact-D characteristic BDE formulation.
Admissibility requires finite supersonic states, a finite ordered prefix,
positive last increments above a scale-aware roundoff guard, the principal
$\beta\in(0,\pi/2)$ branch outside the integrator's tangent fallback, and
a derived angle inside the existing solver bounds. No angle is clipped.
Those bounds are $[-10,65]$ degrees. The positive-increment guard is
$64\epsilon_{\mathrm{mach}}\max(|L|,|R_E|,|R_t|)$ in length units,
and the existing tangent fallback threshold is $10^{-12}$. They are
numerical domain conventions, not additional literature physics.

Invalid trial branches may use an explicitly flagged finite extension only
to evaluate guiding domain penalties and derivatives. A typed final-domain
refusal must prevent their promotion even if the optimizer reports success.
Acceptance requires host/JAX and finite-difference derivative checks,
roundoff endpoint closure, invalid-branch refusal, preserved defaults and
the unchanged independent flow/topology gates.

The focused implementation suite completed with 114 passed and two slow tests
deselected in 57.83 seconds (exit 0). After adding an explicit early nonfinite
refusal and report-only worst-field coordinates, the affected final suite
completed with 42 passed in 19.21 seconds (exit 0). The records are
\code{validation/t18_result.json}, \code{validation/t18_focused_pytest.log} and
\code{validation/t18_final_domain_pytest.log}. The derivative test compares the
entire reduced residual Jacobian with centered finite differences and verifies
the exit-angle dependencies on kdf, exit Mach and prefix state. Separate tests
make both optimizers claim success on invalid boundary data and require refusal.
The tests therefore examine mathematical dependencies and acceptance behavior,
not only successful compilation.

The final independent code review found no material defect within
this scope and confirmed the remaining applicability limits. Actual
profile/refinement acceptance is recorded separately from these component
checks. The benchmark runner now
forwards the opt-in attachment, capacity, seed-accuracy, turning-step and
physics-weight controls without changing shipped case defaults or thresholds.

\subsection{Three solved levels: endpoint fixed, field comparison still fails}
With exact D and E, seed budget 25, requested seed tolerance $10^{-8}$,
physics weight 1 and the other Iteration-13 physical inputs unchanged,
all three independently solved levels pass BVP, all eleven raw physical
gates, Eq.~(5), full-CDE thrust, and export. Each raw exit is
$(x_E,r_E)=(2.94,2.2300224214)R_t$. No endpoint forcing or monotonic
cleanup is used by export.

\begin{longtable}{@{}P{0.23\linewidth}P{0.15\linewidth}P{0.18\linewidth}P{0.19\linewidth}P{0.14\linewidth}@{}}
\toprule\textbf{$n_c/n_k$, turning step} & \textbf{Evaluations} & \textbf{$C_F$} & \textbf{Max weighted residual} & \textbf{C+ max, degrees}\\\midrule
\endhead
16/24, $0.25^\circ$ & 142 & 1.57922759 & $3.375\times10^{-4}$ & 0.0012154\\
24/32, $0.125^\circ$ & 297 & 1.57936995 & $1.455\times10^{-4}$ & 0.0003450\\
32/48, $0.0625^\circ$ & 440 & 1.57948822 & $8.252\times10^{-5}$ & 0.0001469\\
\bottomrule
\end{longtable}

The coarse-to-middle comparison still fails coverage 0.883289, maximum
field angle 0.079660 degrees, and exit-angle drift 0.051755 degrees.
The middle-to-finest comparison passes every metric except maximum
field angle: 0.0637603 degrees against 0.05. Coverage is now 0.917791,
Mach maximum difference 0.0017661, pressure-ratio difference 0.00017654,
and wall difference $0.0004411R_t$. Its worst angle location is
$(x/R_t,r/R_t)=(2.207127684,1.680854790)$: 20.10633858 degrees on the
finest mesh versus 20.17009887 from middle-grid interpolation.
The process correctly exits 2. The exact-endpoint repair is accepted;
that pair's field acceptance is still \emph{FAIL}, as retained in
\code{validation/t18_three_level_exact_exit/report.json}.

\subsection{Actual JAX integration: retain an underresolved result}
A separate real 8/8-point static JAX solve with exact D and E terminates
successfully in 69 LM steps, with admissible boundary geometry and exit
$x_E=2.9399999999999995R_t$. The runtime completes in 9.05 seconds,
but C+ error 0.0097119 degrees and axial mass conservation fail their
independent physical gates. It remains \code{geometric_approximation}
and \code{converged=False}. Thus the JAX wrapper executes the new
parameterization, while correctly refusing to promote an underresolved
field. Runtime success is not recorded as a mesh-validation pass.
See \code{validation/t18_jax_runtime_probe.json}.

\section{Iteration 19: Resolve the remaining field-angle gate}
\subsection{Additional refinement and unchanged limits}
The maximum angle error decreases from 0.07966 to 0.06376 degrees, while
its RMS decreases from 0.02120 to 0.009683 degrees. Other final-pair
metrics pass and the error is localized at a measured physical coordinate.
One further pair was therefore evaluated with the same angle
threshold: the existing 32/48-point, $0.0625^\circ$ profile was compared with
48/64 points, $0.03125^\circ$ and a 2,000-step kernel capacity.
All physical inputs, exact-boundary modes, seed controls, weights,
regularization and acceptance thresholds remain identical. The optimizer
budget remains 1,000 function evaluations. The previous full solution object
was not retained by the three-level process, so the 32/48 baseline is
recomputed before the additional solve; its reported values reproduce the
earlier result.

\subsection{Completed result and independent review}
The additional solve terminates by \code{ftol} after 891 function
evaluations. Its complete kernel stores 1,770 rows against the explicit
2,000-step capacity, with $\theta_B=27.76061475$ degrees. The maximum
weighted residual is $3.8072\times10^{-5}$, C+ maximum
$0.0000427292$ degrees and C- maximum $0.0000112508$ degrees.
All eleven raw physical/topology gates, BVP, Eq.~(5), thrust, admissible
boundary domain and export pass. The raw exit is the requested endpoint,
without post-solve forcing or cleanup. Its $C_F=1.57955115$ is approximately
0.212 percent below the published 1.5829, subject to the explicitly
different throat and starting-line models.

\begin{longtable}{@{}P{0.47\linewidth}P{0.23\linewidth}P{0.20\linewidth}@{}}
\toprule\textbf{Common-coordinate metric} & \textbf{Measured} & \textbf{Unchanged limit}\\\midrule
\endhead
Maximum Mach difference & 0.00080366 & $\leq0.005$\\
Maximum flow-angle difference & 0.0228912 degrees & $\leq0.05$ degrees\\
Maximum pressure-ratio difference & 0.0000749481 & $\leq0.001$\\
Minimum bidirectional field coverage & 0.950270 & $\geq0.90$\\
Maximum wall-radius difference / $R_t$ & 0.000226733 & $\leq0.002$\\
Angle-$B$ difference & 0.0017911 degrees & $\leq0.05$ degrees\\
Exit-angle difference & 0.0090863 degrees & $\leq0.05$ degrees\\
D-fraction difference & 0.0000082903 & $\leq0.0005$\\
Relative thrust-coefficient difference & 0.0000398443 & $\leq0.001$\\
\bottomrule
\end{longtable}

The final pair passes every declared field, wall and global-output criterion,
over 2,546 comparison coordinates. Its maximum angle discrepancy is at
$(2.479149212,1.930288975)R_t$: 18.88536963 degrees on the 32/48 grid
versus 18.90826085 from refined-grid interpolation. Field angle RMS is
0.0039017 degrees and the 99th percentile is 0.014558 degrees.
This is a passing finite-resolution comparison, not an observed-order
proof or a continuum-error bound. No fifth level is needed for the
specified acceptance test.

\buildfig{figures/moc_audit/fixed_context_refinement.pdf}
{Completed exact-boundary, accurate-seed comparisons. Each label gives the
control/starting-point counts of two independent solves; the throat turning
limit halves at each successive level. The earlier pairs retain their failed
angle or coverage results. Only the final pair passes every declared metric.
Connecting the measured points does not establish an asymptotic convergence
order or a continuum-error bound.}{0.98}

The pair runner reports exit 0 after 1,497.54 seconds, including a
282.56-second recomputed baseline and 1,213.29-second refined solve.
This computational cost is material; the accepted profile is not silently
made the default for every caller. The complete record is
\code{validation/t19_refinement/report.json}. Compact field arrays are banked
alongside it; larger local solution pickles remain working artifacts.
The packaged reproduction command is:
\begin{lstlisting}
.venv-jax/bin/python -m scripts.audit_moc_refinement_pair \
  --output-dir builds/moc_refinement_pair_check
\end{lstlisting}
It runs the accepted pair directly, omitting earlier coarse levels. The runner
was packaged from the completed experiment with unchanged configuration and
acceptance expressions; import, direct/module help, compilation and AST
equivalence checks pass. The expensive solve was not rerun merely to change
its file location. See \code{validation/t19_packaging_result.json}.

An independent re-evaluation verified the current source and saved-array
hashes, recomputed bidirectional linear interpolation from the NPZ payloads,
and calculated pressure directly from the constant-$\gamma$ isentropic formula.
Mach, angle, pressure and coverage reproduce the reported metrics to
floating-point tolerance and all pass. This re-evaluation avoids the
production comparison wrapper. Other physical flags were reviewed from
the completed solves, not independently resimulated. The record is
\code{analysis/final_field_reevaluation.json}. The pair satisfies the declared
numerical policy for its documented configuration.

All four accurate-seed searches stop at the angular bracket floor rather
than satisfying the requested $10^{-8}$ seed criterion. Their measured
best errors are explicitly reported, including $2.486\times10^{-8}$
on the added grid. That is not mislabeled convergence: the seed is an
initial frozen-kernel construction, while exact endpoint closure and all
final acceptance tests are evaluated on the solved state.

\subsection{Post-freeze integration checkpoint}
The separate post-freeze suite of point-process failure, characteristic
pairing, residual, stationarity, radial-validity, CLI and source-reference
tests completes with 69 passed in 32.06 seconds (exit 0). It is recorded in
\code{analysis/post_freeze_integration_result.json}. The 114-test component
suite, final 42-test affected suite and this 69-test integration suite have
overlapping coverage and are not summed into a claimed unique-test total.
The older 296-pass/two-failure checkpoint remains explicitly pre-capacity
repair, with its separate source inventory and known LEA/TOP failures.

\section{Final audit and remaining evidence}
\subsection{What is established, and what a green local gate means}
The characteristic-family signs and axisymmetric source term agree with the
independently checked governing equations. The audit found several real
implementation and integration defects, rather than evidence that the entire
MOC formulation was unusable: ignored radius input, accepted exhausted point
iterations, accepted incomplete kernels, inactive optimizer columns, a
misdirected CLI validity check, and incomplete model identity in refinement
comparisons. These repairs have bounded acceptance tests and retained source
evidence. Exact D and E elimination are deliberately opt-in numerical
parameterizations of the prescribed boundary problem.

A successful nonlinear optimizer, a small weighted residual, a locally
compatible characteristic net, a source-contour comparison, an accepted
common-coordinate refinement pair and physical validation are different
claims. The structured report keeps their decisions separate. The numerical
policies in this chapter were not relaxed to make a failed attempt green.
The two preexisting LEA/TOP benchmark status failures remain visible; they
are not evidence that these MOC equations failed their source controls.

\subsection{Research questions requiring additional primary evidence}
The original KL-1969 coefficient table and velocity normalization have not
been established. The removed verification claim does not prove the existing
coefficients wrong; it makes their evidentiary status accurate. Resolving
that question requires an accessible original table or another primary
derivation, followed by independent transonic benchmarks. It would be
incorrect to replace the coefficients solely to improve the current gates.
The exact-axis limit remains a targeted research question: the audit did not
establish an active defect in the current midpoint callers.

\subsection{Implemented capabilities with incomplete validation}
Host/JAX residual and derivative parity verifies the examined discrete
formulation, not convergence of every JAX optimizer solve. The new live-kernel
guard prevents acceptance of an incomplete final march but does not certify
every trial iterate. Reduced-coordinate sensitivity output is explicitly
refused until the complete derivative chain is supported. A documented
single-context passing profile does not define a general accuracy envelope.
Additional independently sourced full-flow reference cases and parameter-domain
studies remain necessary before a broad benchmark-validation label is justified.

\subsection{Architecture and scope boundaries}
The default MDO contour provider remains TOP/B\'ezier; this work does not
silently switch it to a variational MOC inner loop. The general solver defaults
are preserved, and a numerical profile must explicitly request its refinement,
resource and attachment controls. The standard contour conversion can reconcile
the throat join with a cubic bridge; that CAD reconciliation is distinct from
qualifying the raw MOC wall and does not establish physical manufacturability
or measured nozzle performance.

Reacting flow, a thermodynamically consistent variable-gas formulation,
shock fitting, viscous/separation prediction and a general differentiable
forward-wall net are potential extensions outside the current gate-repair
scope. They require their own equations, boundary specifications and independent
validation. The knowledge-graph coverage gap is delivered as an explicit report
and traceable research payload; unrelated graph stores were not bulk-rewritten
or marked reviewed merely because this audit ran.
